% Transcription — fonds Alexandre Grothendieck, shelfmark 14, batch 6 (pages 101-120).
% Established from the facsimile archives/batches/14/batch-06.pdf (a mirror of
% https://grothendieck.umontpellier.fr/14.pdf), per the skill
% transcribe-grothendieck. The batch file carries no cover sheet: PDF page N
% is page 100+N, as the Montpellier footer printed on each PDF page confirms,
% and the archivists' pencil numbers bottom-left agree with that position on
% every leaf (101 to 120 all read).
%
% Pass: Opus 5.5 (claude-opus-5-5), 2026-09-23 — first pass, unchecked against the pages by a human.
%
% Transcribed: 102, 105-108, 110-117, 119, 120 (fifteen pages), plus the cover
% 101 as a section heading. Skipped:
%   101      A cover sheet in his hand (« Conjectures de Lefschetz pour cycles
%            algébriques : Extension aux Jac. inter. / Conjectures locales »).
%            No \page; his words are the section heading.
%   103      A typed leaf of an EGA typescript (its p. « - 115 - »:
%            Corollaires 1-3, formula (13) e + codim >= deg.tr. + codim, (14)),
%            underlined and highlighted by other hands, with corrections of his
%            to the typescript itself (a « Variante » dim O_x <= dim O_y +
%            dim O_x (x) k(y), « et réduits », « cond. des chaînes ») and a few
%            fast lines up the left margin that could not be read. Like the
%            EGA typescript of folder 12 batch 4 and the SGA 3 leaves of this
%            folder's batch 7, left out.
%   104, 109, 118  Blank.
% No letter in this batch (the folder's 1965 letter is not among these pages).
% No date on any leaf; no pagination of his own.
%
% What the batch is. Four runs, all on the Lefschetz-type conjectures for
% groups of cycle classes.
%   102      X a smooth projective surface, Y a smooth hyperplane section:
%            N^1, J^1, A^1 restricted to Y; the kernel K of A^1(X) -> A^1(Y)
%            is of finite type, and « en général » finite; a struck question.
%   105-108  « Conjectures "Lefschetz" locales » (his underlined title,
%            p. 105): S = Spec A local, X = S minus the closed point,
%            Y = X ∩ V(f), U = X - Y. Items (1)-(4): vanishing of N^i, J^i,
%            A^i mod finite groups above the middle, bijectivity/injectivity
%            of J^i(X) -> J^i(Y), N^i(X) -> N^i(Y), A^i(X) -> A^i(Y) below it,
%            « cas critiques »; (3) as a consequence; (4) « Pb »: intersection
%            forms on N^i(X), and vanishing cycles. Worked examples N = 2, 3, 4
%            (Mumford-Hodge for N = 2; « groupes essentiels », « homs
%            essentiels », « Pb » of a natural complement of J^1(X) in J^1(Y),
%            « J^1(X) ≃ J^1(Y)^π », the Hodge index theorem as negativity of a
%            form on N^1(X)), then the cases N = 2m+1, N = 2m.
%   110-115  « Conjectures "Lefschetz" (k alg. clos) » (p. 110): X^n smooth
%            projective, Y^{n-1} a smooth hyperplane section; A^i_r, A^i_a,
%            A^i_τ. (1) restriction A^i(X) -> A^i(Y), (2) Gysin A^r_i(Y) ->
%            A^r_i(X), (3) hard Lefschetz E^{n-2i}, E^{n-2i+1} and the
%            consequence for E^p, (2') the affine complement U, (3') the affine
%            cone E over X, with the primitive parts P_a, I_r. P. 113 is an
%            interpolated « Th. 4 » (cd U ≤ dim A for U = S_f, S = Spec A local
%            noetherian complete), proof begun. Pp. 114-115 check (1), (2),
%            (3), (2') for n = 2, 3, 4 with « ok », « ? » and questions
%            « continue » / « numérique » in the margin.
%   116      « Notations »: H^*, P^*, A^* (Chow ring), N^*, J^*, the exact
%            sequence 0 -> J^* -> A^* -> N^* -> 0, the primitive parts M^*, I^*.
%   117      1) finiteness of « familles algébriques » of numerically trivial
%            cycles, 2) (x,y) - (x,b) - (a,y) + (a,b) ~_a 0, and the consequence:
%            the A^i_r(X)^a are abelian varieties, A_0(X)^a ≃ Alb(X).
%   119-120  X projective, M ⊂ K(X): three equivalent conditions for M to be
%            « limited » (bounded families of locally free sheaves, Hilbert
%            polynomials, quotients of O(-n)^N), proofs of the implications;
%            then (p. 120) for X non-singular and M' ⊂ Z^i(X), five conditions
%            (i)-(v) of boundedness of a family of cycles. The list is closed
%            on p. 122 (batch 7: « Ceci prouve … les implications (ii) ⇒ (i) ⇔
%            (iv) ⇒ (iii) ⇒ (v) »).
%
% Continuations at the batch edges. P. 102 is not continued from p. 100
% (batch 5); the cover 101 opens a run. P. 119 ends mid-sentence (« donc les
% F+R sont quotients d'un … ») and p. 120 takes the argument on; p. 120 ends on
% condition (v), and its conclusion is on p. 122 (batch 7, across the SGA 3
% typescript leaf 121).
%
% Notation fixed for the batch. \mathcal{A} for his cursive A (Chow groups:
% \mathcal{A}^{i}_{r} rational, \mathcal{A}^{i}_{a} algebraic,
% \mathcal{A}^{i}_{\tau} τ-equivalence; lower index \mathcal{A}^{r}_{i} for
% cycles of dimension i). J^{i}, N^{i}, M^{i}, I^{i}, P^{i}, B^{1} upright
% capitals as he writes them; his superscript « a » (the part algebraically
% equivalent to 0) and « τ » kept as exponents. E^{p} and L^{p} as he writes
% them for the Lefschetz operator. « exc. p » (except p-torsion) is kept as
% his abbreviation, in square brackets where he brackets it. « mod gpes
% finis » expanded to « mod groupes finis » only where he writes it out.
% Prose ellipses are « … ».
%
% Legibility. 105-108, 110-112, 114-117 are written fair and read almost
% end to end; 113 is a fast draft whose connecting prose is largely \ill{};
% 119-120 are rough and heavily corrected, and the struck passages are given
% as \struck{\ill{}} where they cannot be read.
\documentclass[11pt,a4paper]{article}
% Le préambule commun à toutes les transcriptions du fonds Grothendieck.
%
% Il tient en une idée : **l'incertitude se compose, elle ne se résout pas.**
% Une transcription qui lisse un mot douteux a détruit la seule chose pour
% laquelle on la faisait — un lecteur doit pouvoir savoir, à chaque endroit,
% ce qui était sur le feuillet et ce qui a été ajouté. Les six macros qui
% suivent sont donc l'appareil critique, et non de la décoration.
%
% Les mêmes commandes sont reconnues par `scripts/render.mjs`, qui produit la
% vue de lecture du site. Une macro ajoutée ici doit l'être aussi là-bas,
% faute de quoi le rendu échouera bruyamment — ce qui est voulu : un rendu
% silencieusement faux est pire qu'un rendu absent.

\ProvidesPackage{grothendieck}[2026/08/08 Transcriptions du fonds Grothendieck]

\usepackage{amsmath,amssymb,amsthm}
\usepackage{mathrsfs}
\usepackage{stmaryrd}
% Les diagrammes commutatifs sont partout dans ce fonds : c'est la langue
% même de la géométrie algébrique de Grothendieck, et une transcription qui
% les rendrait en prose ne transcrirait rien.
\usepackage{tikz-cd}
% Français par défaut — c'est la langue des feuillets — mais l'anglais est
% chargé aussi : la traduction et le résumé sont des documents anglais, et la
% typographie française (espaces avant les deux-points) y serait une faute.
% Ces documents basculent par \selectlanguage{english} en tête de corps.
\usepackage[english,french]{babel}
\usepackage{fontspec}
\usepackage{geometry}
\geometry{a4paper,margin=25mm,marginparwidth=28mm}
\usepackage{marginnote}
\usepackage{xcolor}
% ulem, not soul. `soul`'s \st and \ul are built on a character-by-character
% scan that XeLaTeX's font machinery breaks: the document fails with an
% undefined control sequence rather than degrading. ulem does the same two jobs
% and is Unicode-engine safe. `normalem` stops it hijacking \emph, which the
% transcriptions use for Grothendieck's own emphasis.
\usepackage[normalem]{ulem}
\usepackage{hyperref}
\hypersetup{colorlinks=true,linkcolor=black,urlcolor=black,pdfborder={0 0 0}}

% --- Métadonnées du lot -----------------------------------------------------
% Reprises telles quelles par le script de rendu, qui les affiche en tête de
% la vue de lecture. Elles fixent ce que le fichier prétend couvrir : sans
% elles, une transcription flotte sans point d'attache dans un fonds de seize
% mille feuillets.

\newcommand{\folder}[1]{\gdef\grothFolder{#1}}
\newcommand{\batch}[1]{\gdef\grothBatch{#1}}
\newcommand{\leaves}[2]{\gdef\grothFirst{#1}\gdef\grothLast{#2}}
\newcommand{\foldertitle}[1]{\gdef\grothFoldertitle{#1}}
\gdef\grothFolder{?}\gdef\grothBatch{?}\gdef\grothFirst{?}\gdef\grothLast{?}\gdef\grothFoldertitle{}

% --- Le résumé --------------------------------------------------------------
%
% Il ouvre la lecture modernisée, sous le titre « Résumé », et il est composé
% en petit. Ce n'est pas un ornement : c'est le seul passage du document qui
% ne soit pas des mathématiques, et un lecteur doit pouvoir le reconnaître
% avant de l'avoir lu — puis le sauter s'il connaît déjà le sujet, ou s'y
% arrêter s'il ne le connaît pas. Le filet qui le suit marque la reprise du
% corps.

\newenvironment{resume}
  {\small\setlength{\parskip}{0.45em}}
  {\par\medskip\hrule\medskip}

% --- Mots-clés ---------------------------------------------------------------
%
% En anglais, en clôture du résumé de la lecture modernisée : le vocabulaire
% moderne sous lequel chercher ce que les feuillets construisent. Le manifeste
% du site (`npm run manifest`) les extrait pour étiqueter la cote entière —
% cette ligne est la source unique des tags, il n'y a pas de fichier à côté.
\newcommand{\keywords}[1]{\par\smallskip\noindent
  {\footnotesize\textsc{Keywords} --- \textit{#1}}\par}

% --- L'appareil critique ----------------------------------------------------

% Le numéro de feuillet, en marge. C'est l'ancre qui permet de régler un
% désaccord sur une formule en regardant *un* feuillet plutôt qu'une chemise
% de six cents. Le site s'en sert aussi pour tourner les pages du fac-similé
% quand on fait défiler la transcription.
\newcommand{\leaf}[1]{%
  \marginnote{\footnotesize\color{gray!70}\textsf{#1}}%
  \label{leaf:#1}\ignorespaces}

% Illisible : on ne devine pas. Un mot inventé qui se lit comme les autres est
% le pire résultat possible.
\newcommand{\ill}{\textcolor{red!60!black}{[\,\ldots\,]}}

% Lecture douteuse : le mot est donné, et signalé comme incertain.
\newcommand{\uncertain}[1]{\uline{#1}}

% Ajout de l'éditeur — une lettre manquante, un mot sous-entendu.
\newcommand{\add}[1]{\textcolor{blue!50!black}{[#1]}}

% Biffé par Grothendieck lui-même. Ce qu'il a rayé fait partie du document :
% une rature dit souvent où la pensée a changé de direction.
\newcommand{\struck}[1]{\sout{#1}}

% Note du transcripteur, sur l'état matériel du feuillet.
\newcommand{\note}[1]{\par\smallskip{\small\color{gray!60!black}\textit{[#1]}}\par\smallskip}

% Note marginale de Grothendieck, distincte du corps du texte.
\newcommand{\marginal}[1]{\par\smallskip\noindent\hspace{1em}%
  {\small\color{gray!60!black}#1}\par\smallskip}

% --- Titre ------------------------------------------------------------------

\AtBeginDocument{%
  \begin{center}
    {\large\bfseries Fonds Alexandre Grothendieck --- cote n\textsuperscript{o}~\grothFolder}\\[2pt]
    {\itshape\grothFoldertitle}\\[4pt]
    {\small feuillets \grothFirst--\grothLast\ (lot \grothBatch)}\\[2pt]
    {\footnotesize\color{gray!60!black}Transcription établie d'après le fac-similé numérisé
     par l'Université de Montpellier.\\
     Les lectures incertaines sont soulignées, les illisibles notées [\,\ldots\,],
     les ajouts entre crochets.}
  \end{center}
  \vspace{1em}}

\endinput


\folder{14}
\batch{6}
\pages{101}{120}
\dating{1965-[vers 1971]}
\watermark{Édition de démonstration}
\foldertitle{Théorie des cycles algébriques : conjectures de Hodge, Tate, Lefschetz, Weil : notes manuscrites (s.d.), lettre (1965).}

\begin{document}

\section{Conjectures de Lefschetz pour cycles algébriques}

\note{titre de sa main, sur le feuillet de garde 101, suivi d'une accolade
réunissant deux lignes : « Extension aux Jac. inter. » (\emph{Jac. inter.}
souligné) et « Conjectures locales » (\emph{locales} souligné)}

\page{102}
$X$ surface projective non sing.

$Y$ section hyperplane non sing.

$N^{1}(X) \to N^{1}(Y)$ \quad surjectif mod gpe fini, noyau de rang
\uncertain{$\rho - 1$}

$J^{1}(X) \to J^{1}(Y)$ \quad injectif exc. $p$.

$\mathcal{A}^{1}(X) \to \mathcal{A}^{1}(Y)$ \quad noyau $K$ tel que
$K \cap J^{1}(X)$ fini, \emph{donc} $K$ de type fini.

Soit $B^{1}(X)$ « partie primitive » de $\mathcal{A}^{1}(X)$, formée des
classes de diviseurs de \uncertain{degré} $0$, donc
\[
B^{1}(X) \longrightarrow J^{1}(Y) \qquad [\,B^{1}(X) \text{ contenant } K\,]
\]
\[
0 \to K \xrightarrow[\text{mod gpes finis}]{\text{inj.}} M^{1}(X) \longrightarrow
\underbrace{N^{1}(Y)/\operatorname{Im} N^{1}(X)}_{\text{partie « mobile » de } N^{1}(Y)}
\]
\note{la lecture « mobile » est douteuse}

\struck{Question. Un \ill{} de $M^{1}(X)$ \ill{} \ill{} \ill{} \ill{}
$N^{1}(Y)/N^{1}(X)$ \ill{} \ill{} contenu dans $K$ i.e. dans \ill{}}
\note{passage barré en zigzag ; au-dessous, un petit schéma « $I \to K$ »}

N.B. « En général » $M^{1}(X) \to N^{1}(Y)/N^{1}(X)$ est injectif mod gpes
finis, donc $K$ fini. Je présume qu'il en est ainsi si les sections
\uncertain{hyperplanes} de $X$ des \uncertain{bons} pinceaux
\uncertain{passant par} $Y$ sont des courbes \emph{irréductibles}…
\note{le dernier mot est souligné deux fois}

\section{Conjectures « Lefschetz » locales}

\page{105}
\note{titre souligné, de sa main}

$S = \operatorname{Spec} A$, \quad $A$ anneau str.\ local \uncertain{noeth.},
$f \in \mathfrak{m}_{A}$,
\[
\begin{cases}
X = S - \text{pt fermé} \\
Y = X \cap V(f)
\end{cases}
\]
$X$, $Y$ réguliers, \quad $U = X - Y$.

\textbf{(1)} \quad \struck{\ill{}} $i = \frac{N}{2}$ \quad
$N^{i}(X) \equiv$ \struck{\ill{}} $0$ \quad (mod gpes finis)
\note{le signe entre $i$ et $\frac{N}{2}$ est surchargé ; $=$ est la lecture
finale}

$i \geqslant \frac{N+1}{2}$ \quad $N^{i}(X) \equiv 0$, \quad $J^{i}(X) = 0$
\quad (mod gpes finis)

i.e.\ $\mathcal{A}^{i}(X) \equiv 0$ \quad (mod gpes finis)

\textbf{(2)} \quad $J^{i}(X) \to J^{i}(Y)$
\[
\begin{array}{lll}
\text{bij.\ [exc.\ } p\text{] si} & i \leqslant \frac{N-2}{2} \\[2pt]
\text{inj.\ [exc.\ } p\text{] si} & i \leqslant \frac{N-1}{2} \\[2pt]
\text{nul si} & i \geqslant \frac{N}{2}
\end{array}
\]
[N.B. Si $i \geqslant \frac{N}{2}$, $J^{i}(Y) = 0$ ! d'ailleurs si
$i \geqslant \frac{N+1}{2}$, \ill{} $J^{i}(X) = J^{i}(Y) = 0$]

$N^{i}(X) \to N^{i}(Y)$
\[
\begin{array}{lll}
\text{bij.\ [exc.\ } p\text{] si} & i \leqslant \frac{N-3}{2} \\[2pt]
\text{inj.\ [exc.\ } p\text{] si} & i \leqslant \frac{N-2}{2} \\[2pt]
\text{nul [mod.\ gpes finis] si} & i \geqslant \frac{N-1}{2}
\end{array}
\]
Cas critiques $N^{\frac{N-1}{2}}$, $N^{\frac{N}{2}}$
\note{les deux indices sont soulignés d'un crochet}

[N.B. si $i \geqslant \frac{N-1}{2}$, on a $N^{i}(Y) \equiv 0$ ; si
$i \geqslant \frac{N}{2}$, \ill{} $N^{i}(X) = N^{i}(Y) = 0$ ?]

\page{106}
Cas critiques $N^{\frac{N-2}{2}}$, $N^{\frac{N-1}{2}}$

\note{le $N$ de ces deux symboles est repassé à l'encre épaisse}

$\mathcal{A}^{i}(X) \to \mathcal{A}^{i}(Y)$ \quad bij.\ si
$i \leqslant \frac{N-3}{2}$, \quad inj.\ si $i \leqslant \frac{N-2}{2}$

Cas critiques $i = \frac{N-2}{2}, \frac{N-1}{2}, \frac{N}{2}$

\marginal{conséquence de 1° et 2°}
\textbf{(3)} \quad $\mathcal{A}^{i}(X) \longrightarrow \mathcal{A}^{i}(U)$
\quad \emph{un isom.} mod gpes finis, \uncertain{tjs}\,; \emph{un isom.} si
$i \leqslant \frac{N-1}{2}$ [NB. si \struck{\ill{}} $i \geqslant \frac{N+1}{2}$,
$\mathcal{A}^{i}(X)$ \uncertain{et} $\mathcal{A}^{i}(U)$ \uncertain{sont}
finis !]

$\mathcal{A}^{i}(Y) \longrightarrow \mathcal{A}^{i+1}(X)$ \quad \emph{nul}
\struck{mod} gpes finis, \uncertain{tjs}\,; nul si $i \leqslant \frac{N-1}{2}$
[NB si $i \geqslant \frac{N}{2}$, $\mathcal{A}^{i}(Y)$ \uncertain{est un gpe
fini} !]

\textbf{(4)} \quad \emph{Pb}. Définir des formes « intersections »
$N^{i}(X) \times N^{i}(X) \to \mathbb{Z}$ ??

Définir \struck{par} $J^{\frac{N-1}{2}}(X) \to J^{\frac{N-1}{2}}(Y)$ resp.\
$N^{\frac{N-2}{2}}(X) \to N^{\frac{N-2}{2}}(Y)$ \uncertain{via} théorie des
cycles évanouissants ------

\note{un trait horizontal traverse la page}

\emph{Ex} \quad $N = 2$

\textbf{(1)} \quad $\boxed{N^{1}(X) \text{ fini}}$ : c'est le th.\ de
\struck{\ill{}} Mumford-Hodge.

\textbf{(2)} \quad vide

\textbf{(3)} \quad \struck{$\mathcal{A}^{1}(X)$} $\to$ vide

Groupe fini essentiel \quad $J^{1}(X)$

\page{107}
\emph{Ex} \quad $N = 3$

\textbf{(1)} \quad $\boxed{\mathcal{A}^{2}(X) \text{ fini}}$ \quad en
particulier $J^{2}(X) = 0$

\textbf{(2)} \quad $\boxed{J^{1}(X) \to J^{1}(Y) \quad \text{injectif exc.}}$
\struck{\ill{}}

\textbf{(3)} \quad \struck{$\mathcal{A}^{2}(X)$} contenu dans 1°)

groupes essentiels \quad $J^{1}(X)$, $N^{1}(X)$ \quad définissant
$\mathcal{A}^{1}(X)$

homs essentiels \quad $J^{1}(X) \longrightarrow J^{1}(Y)$, \quad
$N^{1}(X)' \longrightarrow J^{1}(Y)$ \quad $\bigl|$ \quad
$\mathcal{A}^{1}(X)' \longrightarrow J^{1}(Y)$

\marginal{sous-groupe d'indice fini de $N^{1}(X)$}
\note{la glose en marge se rapporte, par un trait, au $N^{1}(X)'$}

\emph{Pb} \quad définir un complémentaire naturel de $J^{1}(X)$ dans
$J^{1}(Y)$ (correspondant aux « cycles évanescents »)

Donner un sens à la formule
\[
J^{1}(X) \xrightarrow{\ \sim\ } J^{1}(Y)^{\pi}
\]
Interpréter le th.\ de \struck{Hodge} l'indice de Hodge pour surfaces
projectives en termes de la négativité d'une certaine forme définie sur
$N^{1}(X)$… ?

\emph{Ex} \quad $N = 4$

\textbf{(1)} \quad
$\begin{cases}
N^{2}(X) \text{ fini} \\
\mathcal{A}^{3}(X) \text{ fini, en particulier } J^{3}(X) = 0
\end{cases}$

\textbf{(2)} \quad $J^{1}(X) \to J^{1}(Y)$ \quad bij exc.\ $p$

$N^{1}(X) \longrightarrow N^{1}(Y)$ \quad inj.\ exc $p$

$\Longrightarrow \mathcal{A}^{1}(X) \to \mathcal{A}^{1}(Y)$ \quad \emph{bij
exc.} $p$

\textbf{(3)} \quad \ill{}
\note{le chiffre est surchargé et rien ne suit sur la page}

\page{108}
groupes essentiels \quad $J^{2}(X)$, $N^{1}(X)$

hom essentiel \quad $N^{1}(X) \longrightarrow N^{1}(Y)$

\emph{Pb} \quad définir un complémentaire naturel de $N^{1}(X)$ dans
$N^{1}(Y)$ (cycles évanouissants). \struck{Donner} Donner un sens à la
formule
\[
N^{1}(X) \xrightarrow{\ \sim\ } N^{1}(Y)^{\pi}
\]
\note{un trait horizontal sépare ce qui suit}

Cas important \quad $N = 2m+1$, on trouve alors
\[
\mathcal{A}^{m+1}(X) \equiv 0 \qquad \bigl(\text{et } \mathcal{A}^{i}(X)
\equiv 0 \text{ si } i \geqslant m+1\bigr)
\]
[Cas $N = 2m$, alors $N^{m}(X) \equiv 0$ \quad i.e.\
$J^{m}(X) \equiv \mathcal{A}^{m}(X)$, et $\mathcal{A}^{m+1}(X) \equiv 0$,
(et $\mathcal{A}^{i}(X) \equiv 0$ si $i \geqslant m+1$)
\note{le crochet ouvert n'est pas refermé}

\section{Conjectures « Lefschetz » (k alg. clos)}

\page{110}
\note{titre de sa main : « Conjectures » souligné, « "Lefschetz" » ajouté
au-dessus, « (k alg. clos) » à la suite}

Soit $X^{n}$ projectif lisse connexe, $Y^{n-1}$ section hyperplane lisse.

--- $\mathcal{A}^{i}_{r}(X)$ … cycles de codimension $i$ mod équivalence
\struck{\ill{}} rationnelle

\quad $\mathcal{A}^{i}_{a}(X)$ \qquad algébrique

\quad $\mathcal{A}^{i}_{\tau}(X)$ \qquad équivalence $\tau$

\textbf{(1)} \quad $\mathcal{A}^{i}_{r}(X) \longrightarrow \mathcal{A}^{i}_{r}(Y)$
\quad \struck{bijectif [exc.\ $p$] si $2i \leqslant n-2$ i.e.\
$i \leqslant \frac{n-1}{2}$ ; injectif exc.\ $p$ si $2i = n-1$ …
i.e.\ $i \leqslant \frac{n-1}{2}$ ; bijectif \ill{}}
\note{un premier tableau encadré, à droite, est biffé en zigzag ; la suite le
reprend ligne par ligne}

$2i \leqslant n-2$ \quad i.e.\ $i < \frac{n-1}{2}$

\qquad bijectif [exc.\ $p$ ?]

$2i = n-1$ \quad i.e.\ $i = \frac{n-1}{2}$

\qquad injectif [exc.\ $p$,]

\qquad induit $\mathcal{A}^{i}_{r}(X)^{\tau} \longrightarrow
\mathcal{A}^{i}_{r}(Y)^{\tau}$ \quad bijectif [exc.\ $p$]

\qquad $\Longrightarrow$ \quad $\mathcal{A}^{i}_{a}(X) \longrightarrow
\mathcal{A}^{i}_{a}(Y)$ \quad bijectif [exc.\ $p$]

\qquad\qquad $\mathcal{A}^{i}_{\tau}(X) \longrightarrow
\mathcal{A}^{i}_{\tau}(Y)$ \quad injectif

$2i = n$ \quad i.e.\ $i = \frac{n}{2}$
\note{$n$ est ici surchargé, dans les deux membres}

\qquad induit $\mathcal{A}^{i}_{r}(X)^{\tau} \longrightarrow
\mathcal{A}^{i}_{r}(Y)^{\tau}$ \quad injectif [exc.\ $p$]

\textbf{(2)} \quad $\mathcal{A}^{r}_{i}(Y) \longrightarrow \mathcal{A}^{r}_{i}(X)$

$2i \leqslant n-3$ \quad i.e.\ $i \leqslant \frac{n-3}{2}$
\struck{\ill{}} $= \frac{n-1}{2} - 1$
\note{le $3$ de $n-3$ est surchargé}

\qquad bijectif [exc.\ $p$ ?]

$2i = n-2$ \quad i.e.\ $i = \frac{n-2}{2}$
\[
\begin{cases}
\text{surjectif [exc.\ } p\text{]} \\
\text{induit } \mathcal{A}^{r}_{i}(Y)^{\tau} \to \mathcal{A}^{r}_{i}(X)^{\tau}
\quad \text{surj.\ [exc.\ } p\text{]} \\
\text{induit } \mathcal{A}^{a}_{i}(Y) \to \mathcal{A}^{a}_{i}(X)
\quad \text{bij.\ [exc.\ } p\text{]}
\end{cases}
\]
\note{les « [exc.\ $p$] » de cette page sont ajoutés au crayon ou d'une encre
plus pâle}

\page{111}
$2i = n-1$ \quad i.e.\ $i = \frac{n-1}{2}$

\qquad induit $\mathcal{A}^{\tau}_{i}(Y) \longrightarrow
\mathcal{A}^{\tau}_{i}(X)$ \quad surj.\ exc.\ $p$

\textbf{(3)} \quad $\mathcal{A}^{i}_{a}(X) \xrightarrow{E^{n-2i}}
\mathcal{A}^{n-i}_{a}(X)$ \quad $(2i \leqslant n)$

\qquad \struck{injectif exc.\ $p$} bijectif mod groupes finis

\qquad \struck{induit $\mathcal{A}^{i}_{\tau}(X) \to
\mathcal{A}^{n-i}_{\tau}(X)$ bijectif mod groupes finis}

\qquad $\mathcal{A}^{i}_{r}(X)^{\tau} \xrightarrow{E^{n-2i+1}}
\mathcal{A}^{n-i+1}_{r}(X)^{\tau}$ \quad $(2i-1 \leqslant n)$

\qquad\qquad bijectif mod groupes finis
\[
\Longrightarrow \qquad \mathcal{A}^{i}_{r}(X) \xrightarrow{E^{p}}
\mathcal{A}^{i+p}_{r}(X)
\]
\qquad a) injectif mod gpes finis si $p \leqslant n-2i$

\qquad b) surjectif mod gpes finis si $p \geqslant n-2i+1$

\note{un trait horizontal ; le reste de la page est blanc}

\page{112}
\textbf{(2')} \quad Posons $U = X - Y$, qui est une variété affine non sing.
\[
\mathcal{A}^{i}_{r}(U) \simeq \mathcal{A}^{i}_{r}(X)/\operatorname{Im}
\mathcal{A}^{i}_{r}(Y), \quad \text{de même pour } \mathcal{A}_{a}(U) \ldots
\]
$2i \leqslant n-2$ \quad i.e.\ $i \leqslant \frac{n}{2} - 1$ \qquad
$\mathcal{A}^{r}_{i}(U) = 0$ \quad (exc.\ $p$ ?) \struck{\ill{}}

$2i = n-1$ \quad i.e.\ $i = \frac{n-1}{2}$ \qquad
$\mathcal{A}^{a}_{i}(U) = 0$ \quad (exc $p$ ?)

\textbf{(3')} \quad Soit $E$ le cône projetant affine associé à $X$, (variété
quasi-affine non sing.\ de dim \uncertain{$2n+1$}).
\note{on attendrait $n+1$ ; nous lisons ce que la page porte}

$2i \geqslant n+2$ \quad i.e.\ $i \geqslant \frac{n}{2} + 1$, \quad
$n - i \leqslant \frac{n}{2} - 1$ \qquad $\mathcal{A}^{i}_{r}(E)$
\struck{$\equiv$} fini, \quad $\mathcal{A}^{i}_{r}(E)^{a} = 0$

$2i \geqslant n+1$ \quad i.e.\ $i \geqslant \frac{n+1}{2}$ \quad i.e.\
$n - i \leqslant \frac{n-1}{2}$ \qquad $\mathcal{A}^{i}_{a}(E)$ fini

\emph{N.B.} On doit avoir
\[
\begin{cases}
\mathcal{A}^{i}_{a}(E) \simeq P^{i}_{a}(X) & \text{(mod gpes finis) si } 2i \leqslant n \\
\mathcal{A}^{i}_{a}(E) \simeq 0 & \text{mod gpes finis si } 2i \geqslant n+1
\end{cases}
\]
\note{entre les deux lignes, une troisième est barrée :
« $\mathcal{A}^{i}_{\tau}(E)^{a} \simeq P^{i}_{\tau}(X)$ mod gpes finis si
$2i \leqslant n+1$ » ; la lettre $P$ est surchargée à l'encre épaisse sur
une première lettre illisible}
\[
\begin{cases}
\mathcal{A}^{i}_{r}(E)^{a} \simeq I^{i}(X) & \text{mod gpes finis si } 2i \leqslant n+1 \\
\mathcal{A}^{i}_{r}(E)^{a} \simeq 0 & \text{mod gpes finis si } 2i \geqslant n+2
\end{cases}
\]
où
\[
\begin{aligned}
P_{a}(X) &\simeq \operatorname{Ker}\bigl(\mathcal{A}^{i}_{a}(X)
\xrightarrow{L^{n-2i}} \mathcal{A}^{n-i}_{a}(X)\bigr)
&& \text{partie « primitive » de la « coh.\ algébrique »} \\
I_{r}(X) &= \operatorname{Ker}\bigl(\mathcal{A}^{i}_{r}(X)^{a}
\xrightarrow{L^{n-2i+1}} \mathcal{A}^{n-i+1}_{r}(X)^{a}\bigr)
&& \text{partie « primitive » de la V.A.\ intermédiaire}
\end{aligned}
\]
\note{les lettres $P$ et $I$ sont surchargées ; les deux gloses sont en marge
droite}

\page{113}
\note{brouillon rapide, traversé de haut en bas par un long trait au crayon ;
la prose de liaison est en grande partie illisible}

\emph{Th.\ 4} \quad Soit $A$ un anneau local noeth.\ complet \uncertain{tel}
que les composantes irréductibles de $S = \operatorname{Spec} A$ soient
« d'égale dim ». Soit $f \in \mathfrak{m}_{A}$, et $U = S_{f}$. Alors
\struck{\ill{}} $\operatorname{cd} U \leqslant \dim A$.

On peut supposer $A$ intègre [normal] complet [de dim $n$], et on est
ramené à prouver $H^{i}(U, \mathbb{Z}/\ell\mathbb{Z}) = 0$ si $i > n$.
$A$ est fini sur un anneau de \struck{séries} formelles
$A_{0} = k[[t_{1}, \ldots, t_{n}]]$, et on \uncertain{peut supposer} de plus
$K/K_{0}$ \emph{séparable}, quitte à remplacer $A$ \ill{} \ill{} \ill{}
$A_{0}$ \ill{} \ill{} \ill{}. Soit $\delta \in \mathfrak{m}A_{0}$ tel que
$S_{\delta}$ \ill{} sur $S_{0\delta}$ \ill{} \ill{} de $K_{0}$ dans $K$
…). On \struck{\ill{}} \ill{} \ill{} $\delta f = t_{n}$ [\ill{} de la dim.\
\ill{} \ill{} \ill{}]. Posons
$A'_{0} = k[[t_{1}, \ldots, t_{n-1}]]\{t_{n}\}$,
$S'_{0} = \operatorname{Spec} A'_{0}$, on \ill{} \ill{} \ill{} $S \mid S_{0}$
\ill{} d'un $\Sigma'$ étale sur $S'_{0}$ \ill{} \ill{} et un \ill{} (\ill{}
\ill{} à normalité)
\[
S \simeq S_{0} \times_{S'_{0}} S', \qquad S' \ [= \operatorname{Spec} A']
\text{ rev.\ fini normal de } S.
\]
D'ailleurs $f$ provient d'un $f' \in A'$.

\marginal{$S \to S_{0} = \operatorname{Spec} k[[t_{1}, \ldots, t_{n}]]
\to S'_{0} = \operatorname{Spec} k[[t_{1}, \ldots, t_{n-1}]]\{t_{n}\}$}
\note{petit schéma vertical en marge gauche, au niveau de $A_{0}$}

\page{114}
\struck{1) \ill{} $n = 2$}

\emph{Ex.} \quad $n = 2$

\textbf{(1)} \quad $\mathcal{A}^{0}_{r}(X) \to \mathcal{A}^{0}_{r}(Y)$ \quad
bijectif exc $p$ \qquad ok

\qquad $\mathcal{A}^{1}_{r}(X)^{\tau} \to \mathcal{A}^{1}_{r}(Y)^{\tau}$ \quad
injectif exc.\ $p$ \qquad ok

\textbf{(2)} \quad $\mathcal{A}^{r}_{0}(Y) \to \mathcal{A}^{r}_{0}(X)$ \quad
surjectif exc.\ $p$ \qquad ok

\qquad $\mathcal{A}^{r}_{0}(Y)^{\tau} \to \mathcal{A}^{r}_{0}(X)^{\tau}$ \quad
---- \qquad ok

\qquad $\mathcal{A}^{a}_{0}(Y) \to \mathcal{A}^{a}_{0}(X)$ \quad bijectif
exc.\ $p$ \qquad ok \qquad (les deux $\simeq \mathbb{Z}$)

\marginal{\uncertain{du moins si} équiv.\ \uncertain{alg.} des 0-cycles
$=$ équiv.\ d'Albanese…}
\note{la note marginale, séparée par un trait vertical, se rapporte aux deux
lignes $\mathcal{A}^{r}_{0}$}

\textbf{(3)} \quad
$\left.\begin{array}{c}
\mathcal{A}^{0}_{a}(X) \xrightarrow{E^{2}} \mathcal{A}^{2}_{a}(X) \\
\wr\!\!\parallel \qquad\qquad \wr\!\!\parallel \\
\mathbb{Z} \qquad\qquad\ \ \mathbb{Z} \\[4pt]
\mathcal{A}^{1}_{a}(X) \xrightarrow{E^{0} = \mathrm{id}} \mathcal{A}^{1}_{a}(X)
\end{array}\right\}$
\quad bijectif mod gpes finis \qquad OK.

\qquad
$\left.\begin{array}{c}
\mathcal{A}^{0}_{r}(X)^{\tau} \longrightarrow \mathcal{A}^{2}_{r}(X) \\
\wr\!\!\parallel \qquad\qquad \wr\!\!\parallel \\
0 \qquad\qquad\ \ 0 \\[4pt]
\mathcal{A}^{1}_{r}(X)^{\tau} \xrightarrow{\ \mathrm{id}\ } \mathcal{A}^{1}_{r}(X)^{\tau}
\end{array}\right\}$
\quad bij.\ mod gpes finis \qquad OK.
\note{à la première ligne du second groupe, l'exposant et l'indice du but se
lisent mal ; on attendrait $\mathcal{A}^{3}_{r}(X)^{\tau}$, nul pour $n = 2$}

\textbf{(2')} \quad $\mathcal{A}^{r}_{0}(U) = 0$

\marginal{O.K. si équivalence \uncertain{numérique} $=$ équiv.\ d'Albanese}

\emph{Ex.} \quad $n = 3$

\textbf{(1)} \quad $\mathcal{A}^{0}_{r}(X) \to \mathcal{A}^{0}_{r}(Y)$ \quad
bij.\ exc $p$ \qquad O.K.

\qquad $\mathcal{A}^{1}_{r}(X) \to \mathcal{A}^{1}_{r}(Y)$ \quad injectif
exc $p$ \qquad O.K.

\qquad $\mathcal{A}^{1}_{r}(X)^{\tau} \to \mathcal{A}^{1}_{r}(Y)^{\tau}$ \quad
bijectif exc $p$ \qquad O.K.
\note{les exposants de cette ligne se lisent mal ; nous suivons le cas
$2i = n-1$ de la page 110}

\textbf{(2)} \quad $\mathcal{A}^{r}_{0}(Y) \to \mathcal{A}^{r}_{0}(X)$ \quad
bijectif exc.\ $p$ \qquad O.K.

\marginal{si équiv.\ \uncertain{num.}\ $=$ équiv.\ Alb.}

\qquad $\mathcal{A}^{a}_{1}(Y) \longrightarrow \mathcal{A}^{a}_{1}(X)$ \quad
surjectif exc.\ $p$ \quad ( ? )
\note{le point d'interrogation est encerclé, et relié par un trait à la
ligne suivante}

\textbf{(3)} \quad
$\begin{array}{c}
\mathcal{A}^{0}_{a}(X) \xrightarrow{E^{3}} \mathcal{A}^{3}_{a}(X) \\
\wr\!\!\parallel\mathbb{Z} \qquad \wr\!\!\parallel\mathbb{Z} \\
\mathcal{A}^{1}_{a}(X) \xrightarrow{E^{1}} \mathcal{A}^{2}_{a}(X)
\end{array}$
\quad \struck{bijectif mod gpes finis O.K.} \struck{bijectif} ( ? )
\quad bijectif mod gpes finis

\marginal{décid.\ \uncertain{gpes finis}}
\note{le second « ( ? ) », cerclé à l'encre épaisse, porte la légende
« Question “numérique” »}

\page{115}
$\mathcal{A}^{0}_{r}(X)^{\tau} \xrightarrow{E^{4}}
\mathcal{A}^{4}_{r}(X)^{\tau}$ \quad bijec.\ mod finis \qquad ok

\qquad\qquad $\wr\!\!\parallel$ \ $0$ \qquad $\wr\!\!\parallel$ \ $0$

\qquad $\mathcal{A}^{1}_{r}(X)^{\tau} \xrightarrow{E^{2}}
\mathcal{A}^{3}_{r}(X)^{\tau}$ \quad id \qquad ok

\qquad $\mathcal{A}^{2}_{r}(X)^{\tau} \xrightarrow{\ \mathrm{id}\ }
\mathcal{A}^{2}_{r}(X)^{\tau}$ \quad id \qquad ok

\marginal{si équiv.\ \ill{} $\simeq$ équiv.\ \ill{}}
\note{la marge se rapporte à la deuxième ligne ; « id » dans la colonne des
conclusions vaut « idem » ; l'exposant $4$ du but est surchargé}

\textbf{(2')} \quad $\mathcal{A}^{r}_{0}(U) = 0$ \qquad
$\mathcal{A}^{a}_{1}(U) = 0$ \quad (exc $p$) \qquad ??
\note{les deux points d'interrogation sont cerclés au crayon}

\emph{Ex.} \quad $n = 4$

\textbf{(1)} \quad $\mathcal{A}^{0}_{r}(X) \to \mathcal{A}^{0}_{r}(Y)$ \quad
bij.\ exc.\ $p$ \qquad ok \qquad (les deux $\simeq \mathbb{Z}$)

\qquad $\mathcal{A}^{1}_{r}(X) \to \mathcal{A}^{1}_{r}(Y)$ \quad bij.\ exc.\
$p$ \qquad O.K.

\qquad \struck{\ill{}}

\qquad $\mathcal{A}^{2}_{r}(X)^{\tau} \to \mathcal{A}^{2}_{r}(Y)^{\tau}$ \quad
injectif exc.\ $p$ \quad ( ? ) --- question « continue »

\textbf{(2)} \quad $\mathcal{A}^{r}_{0}(Y) \to \mathcal{A}^{r}_{0}(X)$ \quad
bijectif exc.\ $p$ \qquad O.K.

\marginal{si équiv.\ \uncertain{num.}\ $=$ équiv.\ Alb}

\qquad $\mathcal{A}^{r}_{1}(Y) \longrightarrow \mathcal{A}^{r}_{1}(X)$ \quad
surj.\ exc.\ $p$ \quad ?

\qquad $\mathcal{A}^{r}_{1}(Y)^{\tau} \longrightarrow
\mathcal{A}^{r}_{1}(X)^{\tau}$ \quad surj.\ exc.\ $p$ \quad ( ? ) --- question
« continue »

\qquad $\mathcal{A}^{a}_{1}(Y) \longrightarrow \mathcal{A}^{a}_{1}(X)$ \quad
bij.\ exc.\ $p$ \quad ( ? ) --- question « numérique » (\uncertain{ess.}\add{t})

\textbf{(3)} \quad $\mathcal{A}^{0}_{a}(X) \xrightarrow{E^{4}}
\mathcal{A}^{4}_{a}(X)$ \quad bij.\ mod gpes finis \qquad O.K.
\qquad (les deux $\simeq \mathbb{Z}$)

\qquad $\mathcal{A}^{1}_{a}(X) \xrightarrow{E^{2}} \mathcal{A}^{3}_{a}(X)$
\quad id \quad ( ? )

\qquad $\mathcal{A}^{2}_{a}(X) \xrightarrow{\ \mathrm{id}\ }
\mathcal{A}^{2}_{a}(X)$ \quad id \qquad O.K.

\note{des traits au crayon relient les « ( ? ) » de cette page à une
mention « mod gpes finis » en marge}

\qquad $\mathcal{A}^{0}_{r}(X)^{\tau} \xrightarrow{E^{5}}
\mathcal{A}^{5}_{r}(X)^{\tau}$ \quad bij.\ id \qquad ok

\qquad\qquad $\wr\!\!\parallel$ \ $0$ \qquad $\wr\!\!\parallel$ \ $0$

\qquad $\mathcal{A}^{1}_{r}(X)^{\tau} \xrightarrow{E^{3}}
\mathcal{A}^{4}_{r}(X)^{\tau}$ \quad id \qquad OK

\marginal{si équiv.\ \uncertain{num.}\ $=$ équiv.\ Alb}

\qquad $\mathcal{A}^{2}_{r}(X)^{\tau} \xrightarrow{\ E\ }
\mathcal{A}^{3}_{r}(X)^{\tau}$ \quad id \quad ( ? ) --- question « continue »

\textbf{(2')} \quad $\mathcal{A}^{r}_{0}(U) = 0$, \quad
$\mathcal{A}^{r}_{1}(U) = 0$ \quad exc $p$. \qquad ??
\note{les deux derniers points d'interrogation, cerclés, sont presque
effacés}

\page{116}
\emph{Notations} \quad Pour $X$ variété projective non singulière.
\[
\begin{cases}
H^{*}(X) & \text{cohomologie } \ell\text{-adique} \\
P^{*}(X) & \text{partie primitive de la cohomologie}
\end{cases}
\]
\marginal{on peut mettre des coeff.}
\[
\begin{cases}
\mathcal{A}^{*}(X) = \mathcal{A}^{*}_{r}(X) & \text{anneau de Chow (équiv.\ rationnelle)} \\
N^{*}(X) = \mathcal{A}^{*}_{a}(X) & \text{id.\ pour équiv.\ algébrique} \\
J^{*}(X) = \mathcal{A}^{*}_{r}(X)^{a}
\end{cases}
\]
\[
0 \longrightarrow J^{*}(X) \longrightarrow \mathcal{A}^{*}(X) \longrightarrow
N^{*}(X) \longrightarrow 0
\]
\[
\begin{cases}
M^{*}(X) & \text{partie primitive de } N^{*}(X) \\
I^{*}(X) & \text{partie primitive de } J^{*}(X)
\end{cases}
\]
\note{« id. » rend ses guillemets de répétition}

\emph{N.B.} Je ne sais pas s'il y a une définition raisonnable de la
« partie primitive » de $\mathcal{A}^{*}(X)$ ---

\page{117}
1) Les cycles num.\ équivalents à $0$ sur $X$, mod équiv.\
\uncertain{rationnelle}, sont décrits par un ens.\ fini de « familles
algébriques » de cycles.
\[
\Longrightarrow
\begin{cases}
(\text{équiv.\ num.} = \tau\text{-équivalence}) \\
\text{les } \mathcal{A}^{i}_{a}(X) \text{ sont de type fini} \ldots
\end{cases}
\ \Downarrow
\]
2) \quad $(x, y) - (x, b) - (a, y) + (a, b) \underset{a}{\sim} 0$

1) et 2) $\Longrightarrow$ Les $\mathcal{A}^{i}_{r}(X)^{a}$ sont des V.A.
\[
\begin{cases}
\mathcal{A}^{i}_{r}(X)^{a} \times \mathcal{A}^{j}_{r}(Y)^{a} \longrightarrow
\mathcal{A}^{i+j}_{r}(X \times Y)^{a} & \text{est une \ill{}} \\
\mathcal{A}^{i}_{r}(X)^{a} \times \mathcal{A}^{j}_{r}(X)^{a} \longrightarrow
\mathcal{A}^{i+j}_{r}(X)^{a} & \text{est une \ill{}}
\end{cases}
\]
\struck{\ill{}} \quad $\mathcal{A}_{0}(X)^{a} \xrightarrow{\ \sim\ }
\operatorname{Alb}(X)$ \quad bijectif
\note{le $\mathcal{A}$ de cette ligne porte un exposant illisible}

2) $\Longleftrightarrow$ \quad $\mathcal{A}_{0}(X \times Y)^{a}
\xrightarrow{\ \sim\ } \mathcal{A}_{0}(X)^{a} \times \mathcal{A}_{0}(Y)^{a}$
\quad bijectif

\page{119}
$X$ schéma projectif

$M \subset K(X)$

\emph{Conditions équivalentes} \quad (i) $\exists$ familles
\emph{limitées} de faisceaux \add{loc.\ libres} $(F_{i})$, $(G_{i})$ sur
$X$, tel que $M$ contenu dans l'ens.\ des $\gamma(F_{i}) - \gamma(G_{i})$
\note{« loc.\ libres » est ajouté au-dessus de la ligne, de sa main}

(ii) $\exists$ des entiers $n$, $N$ et des parties finies
\struck{de polynômes} $P \subset \mathbb{Q}[T]$, telles que les
$\xi \in M$ soient de la forme $\gamma(F) - \gamma(G)$, $F$ et $G$ des
quotients loc.\ libres de $\mathcal{O}(-n)^{N}$, de polynômes de Hilbert
$\in P$.

\marginal{les polynômes de Hilbert des $\xi \in M$ sont en nb fini ?}
\note{ligne cerclée, écrite au-dessus de (iii)}

(iii) [$\exists$ des faisceaux loc.\ libres $A$ et $B$ \ill{}
$\mathcal{O}(-n)^{N}$ \ill{} \ill{}) tels que les
$\xi + \gamma(A)$ \ \uncertain{($\xi \in M$)} soient décrits par des
faisceaux loc.\ libres, quotients de \struck{$B$ = faisceau $\mathcal{O}$
loc.\ libre $(\mathcal{O}(-n)^{N}$ \ill{} \ill{}}].

(iii) $\Longrightarrow$ (ii) trivial, compte tenu que $A$, $B$ sont
\uncertain{quotients} de $\mathcal{O}(-n)^{N}$ pour $n$, $N$ convenables,

(ii) $\Longleftrightarrow$ (i) par la caractérisation des familles limitées

(ii) $\Longrightarrow$ (iii) \quad on a si $G = \mathcal{O}(-n)^{N}/R$, …
\[
\xi = \gamma(F) - \gamma\bigl(\mathcal{O}(-n)^{N}\bigr) + \gamma(R)
\]
d'où
\[
\xi + \gamma\bigl(\mathcal{O}(-n)^{N}\bigr) = \gamma(F + R).
\]
Or $F + R$ est un quotient de $\mathcal{O}(-n)^{N}_{\xi}$, donc
les $F + R$ \uncertain{sont} quotients d'un …
\note{la page s'arrête ici au milieu de la phrase ; l'indice $\xi$ de la
dernière ligne se lit mal}

\page{120}
Supposons \add{de plus} $X$ non singulier, et considérons \struck{des} une
partie $M'$ de $Z^{i}(X)$, \ill{} \uncertain{(considérons les} cycles à
$K$-équivalence près …$)$. Conditions \struck{possibles} :
\note{« de plus » et « considérons les » sont ajoutés au-dessus de la ligne}

(i) \struck{$\exists$} La famille de $M'$ est limitée, \uncertain{i.e.}
définie par une famille limitée de \emph{cycles} algébriques.

(ii) La famille de $M'$ est « numériquement limitée » (i.e.\ un ens.\ fini
de classes d'équivalence numérique \ill{}) et $\exists$
\struck{\ill{}} \add{des}

$Z_{0} \in Z^{i}(X)$ tel que les \struck{$Z + d(Z_{0})$} $Z_{0} + Z$
$(Z \in M)$ soient « positifs » … i.e.\ \uncertain{soient} des cycles
positifs.

(iii) \struck{Il y a une famille limitée de \ill{} de faisceaux loc.\ libres
\ill{} \ill{}} Il y a une famille limitée $M \subset K(X)$ telle que
$\forall Z \in M'$, $\exists \xi \in M$ \uncertain{tel que}
$(i-1)!\, Z = c^{i}(\xi)$ \struck{\ill{}}

(iv) Il y a une famille limitée $M \subset K^{(i)}(X)$ telle que
$M' \subset \operatorname{Image}$ de $M$.

(v) La famille $M'$ est \uncertain{numériquement} limitée.

\note{la suite de cette liste — les implications entre (i) et (v) — est à la
page 122 ; la page 121 est un feuillet dactylographié étranger}

\end{document}
